% Part: computability
% Chapter: computability-theory
% Section: reducibility

\documentclass[../../../include/open-logic-section]{subfiles}

\begin{document}

\olfileid[hi]{cmp}{thy}{red}
\olsection{अपचयनीयता}

\begin{explain}
अब हम जानते हैं कि कम-से-कम एक समुच्चय $K_0$ ऐसा है जो संगणनीयतः
परिगणनीय तो है, पर संगणनीय नहीं। स्पष्टतः ऐसे और भी समुच्चय हैं।
अपचयनीयता की विधि अन्य समुच्चयों में ये गुण दिखाने का एक शक्तिशाली ढंग
देती है, जिसमें बार-बार मूल सिद्धांतों पर लौटने की आवश्यकता नहीं पड़ती।

सामान्यतः, किसी समुच्चय $A$ का किसी समुच्चय~$B$ में ``अपचयन'' वह विधि
है जो प्रत्येक !!{element} के~$B$ में होने या न होने के उत्तर को उसी
!!{element} के~$A$ में होने या न होने के उत्तर में बदलती है। हम
``अनेक-एक अपचयनीयता'' नामक धारणा पर ध्यान देंगे, यद्यपि विभिन्न गुणों
वाली अपचयनीयता की कई अन्य धारणाएँ भी हैं। संगणनात्मक जटिलता के अध्ययन
में भी अपचयनीयता की धारणाएँ केंद्रीय हैं, जहाँ दक्षता के प्रश्नों पर भी
विचार करना पड़ता है। उदाहरणतः, कोई समुच्चय ``NP-पूर्ण'' कहलाता है यदि वह
NP में हो और प्रत्येक NP समस्या को उसमें अपचयित किया जा सके; इसमें नीचे
वर्णित धारणा जैसी अपचयन-धारणा प्रयुक्त होती है, साथ में यह अतिरिक्त
आवश्यकता होती है कि अपचयन बहुपद समय में संगणित हो सके।

हम अपचयन की धारणा का परोक्ष रूप से पहले ही उपयोग कर चुके हैं। समुच्चय~$K$
को परिभाषित करें:
\[
K = \Setabs{x}{\cfind{x}(x) \fdefined},
\]
अर्थात् $K = \Setabs{x}{x \in W_x}$। विराम समस्या के असमाधेय होने का
हमारा प्रमाण (\olref[hlt]{thm:halting-problem}) सबसे सीधे ढंग से दिखाता
है कि $K$ संगणनीय नहीं है। याद करें कि $K_0$ यह समुच्चय है:
\[
K_0 = \Setabs{\tuple{e, x}}{\cfind{e}(x) \fdefined },
\]
अर्थात् $K_0 = \Setabs{\tuple{x,e}}{x \in W_e}$।~$K$ की असंगणनीयता के
किसी भी प्रमाण को~$K_0$ की असंगणनीयता तक बढ़ाना सरल है: यदि $K_0$
संगणनीय होता, तो हम केवल यह पूछकर तय कर सकते थे कि !!a{element}~$x$,
$K$ में है या नहीं कि युग्म $\tuple{x, x}$, $K_0$ में है या नहीं। फलन~$f$,
जो $x$ को $\tuple{x, x}$ पर प्रतिचित्रित करता है, $K$ से $K_0$ में
एक \emph{अपचयन} का उदाहरण है।
\end{explain}

\begin{defn}
$A$ और $B$ को प्राकृतिक संख्याओं के समुच्चय मानें। कोई संगणनीय
फलन~$f\colon \Nat \to \Nat$, $A$ से~$B$ में \emph{अनेक-एक अपचयन} तभी
और केवल तभी है जब प्रत्येक प्राकृतिक संख्या~$x$ के लिए
\[
x \in A \quad \text{तभी और केवल तभी} \quad f(x) \in B.
\]
यदि ऐसा अपचयन $f$ अस्तित्व में हो, तो हम कहते हैं कि $A$,~$B$ में
\emph{अनेक-एक अपचयनीय} है और $A \leq_m B$ लिखते हैं। यदि $A$, $B$ में
अनेक-एक अपचयनीय हो और विलोमतः भी ऐसा हो, तो $A$ और $B$ \emph{अनेक-एक
समतुल्य} कहलाते हैं और $A \equiv_m B$ लिखा जाता है।
\end{defn}

यदि ऊपर की परिभाषा का फलन $f$ एकैकी हो, तो $A$,~$B$ में \emph{एक-एक
अपचयनीय} कहलाता है। नीचे वर्णित अधिकांश अपचयन इस अधिक प्रबल आवश्यकता
को पूरा करते हैं, पर हम इस तथ्य का उपयोग नहीं करेंगे।

\begin{digress}
यह सत्य है, किंतु तनिक भी स्पष्ट नहीं, कि एक-एक अपचयनीयता वास्तव में
अनेक-एक अपचयनीयता से अधिक प्रबल आवश्यकता है। दूसरे शब्दों में, ऐसे अनंत
समुच्चय $A$ और~$B$ हैं जिनमें $A$,~$B$ में अनेक-एक अपचयनीय है, किंतु
~$B$ में एक-एक अपचयनीय नहीं।
\end{digress}

\end{document}
